% Problem record op_845158213592821f. This TeX file is derived from
% database/problems_json/op_845158213592821f.json by scripts/sync-tex.mjs; edit the JSON record
% and rerun the script rather than editing this file.
\section{Unknown-structure Hamiltonian learning from Gibbs states at all temperatures}
\label{sec:op_845158213592821f}

\paragraph{Problem.}
Can every bounded-degree local Hamiltonian with unknown interaction support be learned efficiently from copies of its Gibbs state at any fixed inverse temperature?

Fix constants $k$ and $g$. Let $\mathcal P_{n,k}$ be the nonidentity $n$-qubit Pauli strings of weight at most $k$, and consider

\begin{equation}
H=\sum_{P\in\mathcal P_{n,k}}a_PP,
\quad |a_P|\leq1,
\quad \max_j|\{P:a_P\neq0,\ j\in\operatorname{supp}(P)\}|\leq g,
\qquad
\rho_\beta=\frac{e^{-\beta H}}{\operatorname{Tr}(e^{-\beta H})}.
\label{eq:8451-1}
\end{equation}

The nonzero Pauli terms are unknown, while $\beta>0$ is known and fixed independently of $n$. Can independent copies of $\rho_\beta$ in Eq.~\eqref{eq:8451-1} be used to output coefficients satisfying $\max_P|\widehat a_P-a_P|\leq\varepsilon$ with probability at least $2/3$, using $\operatorname{poly}(n,1/\varepsilon)$ copies and classical time for every fixed $\beta$?

\subsection*{Status}

Unsolved

\subsection*{Source}

This precise formulation is editor wording based on the unresolved direction and limitations documented in the cited primary literature \sourcecite{ref:8451-narayanan25}{Narayanan25}\sourcecite{ref:8451-lewis26}{Lewis26}; it is not presented as a verbatim conjecture of those authors.

\subsection*{Progress}

\begin{itemize}
  \item With a supplied bounded-degree list of $r$ Pauli terms, learning at every fixed temperature is possible. Narayanan's Theorem 1.6 gives
  \begin{equation}
N=O\!\left(r^6(1/\varepsilon)^{O(\beta^2)}
  +\frac{\log r}{\beta^2\varepsilon^2}\right)
\label{eq:8451-3}
\end{equation}
  copies and polynomial computational time for fixed locality, interaction degree, and $\beta$. The supplied list is essential to this theorem; it is parameter learning rather than unknown-structure learning. \sourcecite{ref:8451-narayanan25}{Narayanan25}

The displayed definitions, constraints, and target bounds are recorded in Eqs.~\eqref{eq:8451-3}.

  \item For geometrically local Hamiltonians with a known interaction dictionary, Chen, Anshu, and Nguyen obtain the sharper lattice sample bound
  \begin{equation}
N=\widetilde O\!\left(\frac{e^{\operatorname{poly}(\beta)}}{\beta^2\varepsilon^2}\right)\log(n/\delta),
\label{eq:8451-4}
\end{equation}
  where $\delta$ is the failure probability and the tilde suppresses logarithmic factors. Their all-temperature results do not remove the supplied-structure assumption. \sourcecite{ref:8451-chen25}{Chen25}

The displayed definitions, constraints, and target bounds are recorded in Eqs.~\eqref{eq:8451-4}.

  \item Unknown-structure learning is now solved at sufficiently high temperature. Lewis, Tang, and Wright's Theorem 5.18 proves, under the normalization above,
  \begin{equation}
\beta\leq\frac{1}{1000e^6(2kg+1)^8}
  \quad\Longrightarrow\quad
  N=O\!\left(\frac{\log(n/\delta)}{\beta^2\varepsilon^2}\right),
\label{eq:8451-5}
\end{equation}
  with classical runtime $O(n^k\operatorname{poly}(g)\log(n/\delta)/(\beta^2\varepsilon^2))$. Thus neither high-temperature structure learning nor all-temperature known-structure learning should be listed as open. \sourcecite{ref:8451-lewis26}{Lewis26}

The displayed definitions, constraints, and target bounds are recorded in Eqs.~\eqref{eq:8451-5}.

  \item Section 1.4 of the 29 June 2026 preprint explicitly leaves all-temperature structure learning open. Supplying all $O(n^k)$ candidate Pauli terms to \sourcecite{ref:8451-narayanan25}{Narayanan25} does not immediately solve the problem: their candidate interaction graph no longer has bounded degree. \sourcecite{ref:8451-narayanan25}{Narayanan25}\sourcecite{ref:8451-lewis26}{Lewis26}
\end{itemize}

\subsection*{References}

\begin{enumerate}[label={},leftmargin=4.8em,labelsep=0.5em]
  \footnotesize
  \item[\textup{[Narayanan25]}]\label{ref:8451-narayanan25}
  S. Narayanan, "Improved algorithms for learning quantum Hamiltonians, via flat polynomials," in \emph{Proceedings of the Thirty Eighth Conference on Learning Theory}, Proceedings of Machine Learning Research 291, 4360–4385 (2025). \href{https://proceedings.mlr.press/v291/narayanan25a.html}{Proceedings}; \href{https://arxiv.org/abs/2407.04540}{arXiv:2407.04540}.

  \item[\textup{[Chen25]}]\label{ref:8451-chen25}
  C.-F. Chen, A. Anshu, and Q. T. Nguyen, "Learning quantum Gibbs states locally and efficiently," arXiv preprint (2025), version 1, 3 April 2025. \href{https://arxiv.org/abs/2504.02706}{arXiv:2504.02706}.

  \item[\textup{[Lewis26]}]\label{ref:8451-lewis26}
  L. Lewis, E. Tang, and J. Wright, "Learning the structure of open quantum systems," arXiv preprint (2026), version 1, 29 June 2026. \href{https://arxiv.org/abs/2606.30358}{arXiv:2606.30358}.
\end{enumerate}

\subsection*{Comment}

The unresolved conjunction is unknown interaction support, arbitrary fixed positive temperature, and polynomial resources. Constants and polynomial exponents may depend on $k$, $g$, and $\beta$; this question does not demand efficient scaling as the temperature approaches zero with system size. Lewis, Tang, and Wright explicitly leave all-temperature Gibbs-state structure learning open in Section 1.4; no impossibility result is asserted here.

\subsection*{Field}

Quantum metrology; Quantum algorithm

\subsection*{Topic}

Quantum estimation; Computational complexity and computability; Quantum thermodynamics

\subsection*{ID}

\texttt{op\_845158213592821f}
